\SourceLeaf{82}%
\[
1+d+E+367(A-1700)-2b
\]
or
\[
d+E+367(A-1700)-2b
\]
I then subtract $364(A-1700)$, which gives
\begingroup\postdisplaypenalty=10000
\[
d+E+3A-5099-2b
\]
\endgroup
or
\[
d+E+3A-5100-2b
\]
Adding 5096 next gives
\[
d+E+3A-3-2b
\]
or
\[
d+E+3A-4-2b
\]
Finally, subtracting $3A-3c$, which is clearly divisible by 7, gives
\[
d+E+3c-3-2b
\]
or
\[
d+E+3c-4-2b
\]
Subtracting this from $7c+7d$ gives
\[
3+2b+4c+6d-E
\]
or
\[
4+2b+4c+6d-E
\]
which must therefore be divisible by 7. From this it is clear that $E$ is the remainder obtained by dividing
\[
3+2b+4c+6d
\]
or
\[
4+2b+4c+6d
\]
by 7. Consequently $E=e$.

Easter therefore falls on the $22+d+e^{\mathrm{th}}$ of March or, which is the same thing, on the $d+e-9^{\mathrm{th}}$ of April. W. Z. B. W.

\SourceLeaf{83}%
\begin{center}\textit{Fully General Rules for Calculating Easter\\
in Both the Julian and the Gregorian Calendar.}\end{center}
\begin{center}
\begin{tabular}{l|r|l}
Dividing & by & gives the remainder\\\hline
the year number & 19 & $a$\\
the year number & 4 & $b$\\
the year number & 7 & $c$\\
the number $19a+M$ & 30 & $d$\\
the number $2b+4c+6d+N$ & 7 & $e$\\\hline
\end{tabular}\par
Easter then falls on the $22+d+e^{\mathrm{th}}$ of March\par
or on the $d+e-9^{\mathrm{th}}$ of April
\end{center}

$M$ and $N$ are numbers whose values are constant for all time in the Julian calendar and for at least 100 years at a time in the Gregorian calendar. In the former, $M=15$ and $N=6$. In the latter, from its introduction through 1699, $M=22$ and $N=2$.

\begin{center}\small
\setlength{\tabcolsep}{5pt}\begin{tabular}{rrll|rrll}
from 1700 & to 1799 & $M=23$, & $N=3$ & from 2100 & to 2199 & $M=24$, & $N=6$\\
1800 & $\ldots$1899 & $M=23$, & $N=4$ & 2200 & $\ldots$2299 & $M=25$, & $N=0$\\
1900 & $\ldots$1999 & $M=24$, & $N=5$ & 2300 & $\ldots$2399 & $M=26$, & $N=1$\\
2000 & $\ldots$2099 & $M=24$, & $N=5$ & 2400 & $\ldots$2499 & $M=25$, & $N=1$
\end{tabular}
\end{center}

In general, the values of $M$ and $N$ in the Gregorian calendar for any given century from $100k$ to $100k+99$ are found by the following rule.

Let
\[
k\text{ divided by }\left\{\begin{smallmatrix}3\\4\end{smallmatrix}\right\}
\text{ give the (integer) quotients }
\left\{\begin{smallmatrix}p\\q\end{smallmatrix}\right\}
\]
with the remainders disregarded.
Then
\[
\left\{\begin{smallmatrix}M\\N\end{smallmatrix}\right\}
\text{ is the }\text{\textit{remainder}}\text{ obtained by dividing }
\left\{\begin{smallmatrix}15+k-p-q\\4+k-q\end{smallmatrix}\right\}
\text{ by }\left\{\begin{smallmatrix}30\\7\end{smallmatrix}\right\}\text{.}
\]
\textit{Example.} For the 100 years from 4700 to 4799, $k=47$, $p=15$, $q=11$. Thus $15+k-p-q=36$ and $4+k-q=40$, so $M=6$ and $N=5$. For the year 4763, for instance, this gives
\SourceLeaf{84}%
\begin{center}\small
\begingroup\setlength{\tabcolsep}{5pt}
\begin{tabular}{r|r|l}
$a=13$ & $19a+M=253$ & $e=3$\\
$b=3$ & $d=13$ & Easter on $13+3-9$, that is, the 7th of April\\
$c=3$ & $2b+4c+6d+N=101$ & \quad in the Greg. calendar\\
\multicolumn{3}{@{}c@{}}{}\\[-4pt]
\multicolumn{2}{@{}l|}{In the Julian calendar, however} & \\
\multicolumn{2}{r|}{$19a+M=262$} & $e=2$\\
\multicolumn{2}{r|}{$d=22$} & Easter on $22+2-9$, that is, the\\
\multicolumn{2}{r|}{$2b+4c+6d+N=156$} & \quad 15th of April
\end{tabular}\endgroup
\end{center}

In the \textit{Gregorian calendar} the only exceptions to the above rules are the following two.

I. If the calculation gives the 26th of April for Easter, the 19th of April is \textit{always} taken instead. (E.g. 1609. 1989).

It is easy to see that this case can occur only when the calculation gives $d=29$ and $e=6$. The value of $d$ can be 29 only when dividing $11M+11$ by 30 gives a remainder \textit{less} than 19. For this to happen, $M$ must have one of the following 19 values.
\[
\begin{gathered}
0,\ 2,\ 3,\ 5,\ 6,\ 8,\ 10,\ 11,\ 13,\ 14,\ 16,\ 17,\ 19,\ 21,\ 22,\ 24,\ 25,\ 27,\ 29
\end{gathered}
\]
\Needspace{7\baselineskip}
II. If the calculation gives $d=28$, $e=6$, and if dividing $11M+11$ by 30 also gives a remainder less than 19, then the calculation gives the 25th of April, but Easter falls on the 18th of April. --- It is easy to see that this case can occur only in centuries in which $M$ has one of the following eight values.
\[
2,\ 5,\ 10,\ 13,\ 16,\ 21,\ 24,\ 29
\]
\textit{Apart from these two exceptions, the above rules are completely general.}
\par\Needspace{9\baselineskip}\bigskip\centerline{\rule{.24\textwidth}{.3pt}}\medskip
[Handwritten remark:]

$p$ is determined as the quotient when $8k+13$ is divided by 25.
\begin{center}\small
\begin{tabular}{lrl}
From Easter to Michaelmas there are on average $\ .\ .$ & $173\frac{52}{95}$ & days\\
and (including Easter, excluding Michaelmas) $\ .\ .\ .$ & $25\frac{21}{95}$ & Sundays\\
& $148\frac{31}{95}$ & weekdays\\
From Michaelmas to Easter, however $\ .\ .\ .\ .\ .\ .$ & $164\frac{773}{950}$ & weekdays
\end{tabular}
\end{center}
\par\bigskip\centerline{\rule{.24\textwidth}{.3pt}}
